% --- Archon physics formalization source begin ---
% archon:physics
% archon:covers PhyXMiniProblems/problem_phyx_mini_0306.lean
% archon:source-report reports/phyx_mini/problem_phyx_mini_0306.source.json

\chapter{Physics problem phyx\_mini\_0306}
\label{ch:PhyXMiniProblems_problem_phyx_mini_0306}

\paragraph{Problem source.}
\#\# Physical scenario

In figure, a circular loop of string is set spinning about the center point in a place with negligible gravity. The radius is \$4.00\$ cm and the tangential speed of a string segment is \$5.00\$ cm/s. The string is plucked.

\#\# Question

At what speed do transverse waves move along the string?

\#\# Answer choices

- A: A: 4.4 cm/s

- B: B: 4.6 cm/s

- C: C: 4.8 cm/s

- D: D: 5.0 cm/s

\#\# Figure caption (auxiliary; use the image as primary evidence)

The image shows a simple black circle with a thin outline. To the right of the circle, there is a curved arrow pointing clockwise. The arrow's tail starts near the top right edge of the circle and curves downwards to point outward, suggesting a clockwise direction around the circle. There are no additional objects or text within the image.

\#\# Dataset metadata

Wave Properties; Physical Model Grounding Reasoning, Multi-Formula Reasoning

\paragraph{Recorded answer/context.}
D: D: 5.0 cm/s

\paragraph{Figure/image path.}
/root/proposal\_for\_physic/science-mango/phyx\_mini\_run/phyx\_data/test\_image/306.png

\paragraph{Formalization target.}
create a compiling Lean file with sorry bodies at `PhyXMiniProblems/problem\_phyx\_mini\_0306.lean`. The Lean declarations must preserve the physical quantities, dimensions or dimensional roles, figure labels, governing-law hypotheses, and final relation expressed by this problem.
Use Mathlib/Physlib names found through LeanExplore where available. If a physics API is missing, introduce faithful local abstractions rather than scalar placeholder aliases.

\begin{theorem}[Physics formalization target]
\label{thm:physics:phyx_mini_0306:target}
\lean{PhyXMiniProblems.ProblemPhyXMini0306.problem_phyx_mini_0306}
\uses{def:physics:phyx-mini-0306:phyxminiproblems-problemphyxmini0306-speedincentimeterspersecond, def:physics:phyx-mini-0306:phyxminiproblems-problemphyxmini0306-rotatingcircularstringsetup, def:physics:phyx-mini-0306:phyxminiproblems-problemphyxmini0306-matchesproblemscenario, def:physics:phyx-mini-0306:phyxminiproblems-problemphyxmini0306-matchessuppliedfigure, def:physics:phyx-mini-0306:phyxminiproblems-problemphyxmini0306-matchesproblemreadouts, def:physics:phyx-mini-0306:phyxminiproblems-problemphyxmini0306-hasphysicalparameters, def:physics:phyx-mini-0306:phyxminiproblems-problemphyxmini0306-satisfiesrotatingloopdynamics, def:physics:phyx-mini-0306:phyxminiproblems-problemphyxmini0306-satisfiestransversestringwavelaw, def:physics:phyx-mini-0306:phyxminiproblems-problemphyxmini0306-answerchoice, def:physics:phyx-mini-0306:phyxminiproblems-problemphyxmini0306-answerchoice-speedcentimeterspersecond, def:physics:phyx-mini-0306:phyxminiproblems-problemphyxmini0306-recordeddatasetanswer}
The assigned autoformalize agent should translate this physics problem into Lean declarations in the covered file, with theorem and lemma proof bodies written as `by sorry`.
\end{theorem}
\begin{proof}
This is an autoformalization task, not a proof task. Produce statements that can later be proved by the physics prover without weakening the physical model.
\end{proof}

% --- Archon declaration topology begin ---
\section{Declaration topology}
\begin{definition}[Length Quantity]
  \label{def:physics:phyx-mini-0306:phyxminiproblems-problemphyxmini0306-lengthquantity}
  \lean{PhyXMiniProblems.ProblemPhyXMini0306.LengthQuantity}
  A nonnegative, unit-independent physical length
\end{definition}
\begin{proof}
  This is the declared model definition of Length Quantity.
\end{proof}
\begin{definition}[Mass Quantity]
  \label{def:physics:phyx-mini-0306:phyxminiproblems-problemphyxmini0306-massquantity}
  \lean{PhyXMiniProblems.ProblemPhyXMini0306.MassQuantity}
  A nonnegative, unit-independent physical mass
\end{definition}
\begin{proof}
  This is the declared model definition of Mass Quantity.
\end{proof}
\begin{definition}[Linear Mass Density Quantity]
  \label{def:physics:phyx-mini-0306:phyxminiproblems-problemphyxmini0306-linearmassdensityquantity}
  \lean{PhyXMiniProblems.ProblemPhyXMini0306.LinearMassDensityQuantity}
  Mass per unit length of the uniform string
\end{definition}
\begin{proof}
  This is the declared model definition of Linear Mass Density Quantity.
\end{proof}
\begin{definition}[Acceleration Quantity]
  \label{def:physics:phyx-mini-0306:phyxminiproblems-problemphyxmini0306-accelerationquantity}
  \lean{PhyXMiniProblems.ProblemPhyXMini0306.AccelerationQuantity}
  A nonnegative physical acceleration
\end{definition}
\begin{proof}
  This is the declared model definition of Acceleration Quantity.
\end{proof}
\begin{definition}[Tension Quantity]
  \label{def:physics:phyx-mini-0306:phyxminiproblems-problemphyxmini0306-tensionquantity}
  \lean{PhyXMiniProblems.ProblemPhyXMini0306.TensionQuantity}
  A tensile or resultant force, with dimension mass * length / time²
\end{definition}
\begin{proof}
  This is the declared model definition of Tension Quantity.
\end{proof}
\begin{definition}[Speed Quantity]
  \label{def:physics:phyx-mini-0306:phyxminiproblems-problemphyxmini0306-speedquantity}
  \lean{PhyXMiniProblems.ProblemPhyXMini0306.SpeedQuantity}
  A nonnegative physical speed
\end{definition}
\begin{proof}
  This is the declared model definition of Speed Quantity.
\end{proof}
\begin{definition}[length Readout]
  \label{def:physics:phyx-mini-0306:phyxminiproblems-problemphyxmini0306-lengthreadout}
  \lean{PhyXMiniProblems.ProblemPhyXMini0306.lengthReadout}
  \uses{def:physics:phyx-mini-0306:phyxminiproblems-problemphyxmini0306-lengthquantity}
  Read a length in a selected length unit
\end{definition}
\begin{proof}
  This is the declared model definition of length Readout.
\end{proof}
\begin{definition}[mass Readout]
  \label{def:physics:phyx-mini-0306:phyxminiproblems-problemphyxmini0306-massreadout}
  \lean{PhyXMiniProblems.ProblemPhyXMini0306.massReadout}
  \uses{def:physics:phyx-mini-0306:phyxminiproblems-problemphyxmini0306-massquantity}
  Read a mass in a selected mass unit
\end{definition}
\begin{proof}
  This is the declared model definition of mass Readout.
\end{proof}
\begin{definition}[linear Mass Density Readout]
  \label{def:physics:phyx-mini-0306:phyxminiproblems-problemphyxmini0306-linearmassdensityreadout}
  \lean{PhyXMiniProblems.ProblemPhyXMini0306.linearMassDensityReadout}
  \uses{def:physics:phyx-mini-0306:phyxminiproblems-problemphyxmini0306-linearmassdensityquantity}
  Read a linear mass density in selected mass and length units
\end{definition}
\begin{proof}
  This is the declared model definition of linear Mass Density Readout.
\end{proof}
\begin{definition}[acceleration Readout]
  \label{def:physics:phyx-mini-0306:phyxminiproblems-problemphyxmini0306-accelerationreadout}
  \lean{PhyXMiniProblems.ProblemPhyXMini0306.accelerationReadout}
  \uses{def:physics:phyx-mini-0306:phyxminiproblems-problemphyxmini0306-accelerationquantity}
  Read an acceleration in selected length and time units
\end{definition}
\begin{proof}
  This is the declared model definition of acceleration Readout.
\end{proof}
\begin{definition}[tension Readout]
  \label{def:physics:phyx-mini-0306:phyxminiproblems-problemphyxmini0306-tensionreadout}
  \lean{PhyXMiniProblems.ProblemPhyXMini0306.tensionReadout}
  \uses{def:physics:phyx-mini-0306:phyxminiproblems-problemphyxmini0306-tensionquantity}
  Read tension in the force unit induced by selected base units
\end{definition}
\begin{proof}
  This is the declared model definition of tension Readout.
\end{proof}
\begin{definition}[speed Readout]
  \label{def:physics:phyx-mini-0306:phyxminiproblems-problemphyxmini0306-speedreadout}
  \lean{PhyXMiniProblems.ProblemPhyXMini0306.speedReadout}
  \uses{def:physics:phyx-mini-0306:phyxminiproblems-problemphyxmini0306-speedquantity}
  Read a speed in selected length and time units
\end{definition}
\begin{proof}
  This is the declared model definition of speed Readout.
\end{proof}
\begin{definition}[length In Centimeters]
  \label{def:physics:phyx-mini-0306:phyxminiproblems-problemphyxmini0306-lengthincentimeters}
  \lean{PhyXMiniProblems.ProblemPhyXMini0306.lengthInCentimeters}
  \uses{def:physics:phyx-mini-0306:phyxminiproblems-problemphyxmini0306-lengthquantity, def:physics:phyx-mini-0306:phyxminiproblems-problemphyxmini0306-lengthreadout}
  Centimeter readout used for the stated radius
\end{definition}
\begin{proof}
  This is the declared model definition of length In Centimeters.
\end{proof}
\begin{definition}[speed In Centimeters Per Second]
  \label{def:physics:phyx-mini-0306:phyxminiproblems-problemphyxmini0306-speedincentimeterspersecond}
  \lean{PhyXMiniProblems.ProblemPhyXMini0306.speedInCentimetersPerSecond}
  \uses{def:physics:phyx-mini-0306:phyxminiproblems-problemphyxmini0306-speedquantity, def:physics:phyx-mini-0306:phyxminiproblems-problemphyxmini0306-speedreadout}
  Centimeters-per-second readout used for the given and requested speeds
\end{definition}
\begin{proof}
  This is the declared model definition of speed In Centimeters Per Second.
\end{proof}
\begin{definition}[Rotation Sense]
  \label{def:physics:phyx-mini-0306:phyxminiproblems-problemphyxmini0306-rotationsense}
  \lean{PhyXMiniProblems.ProblemPhyXMini0306.RotationSense}
  Sense of rotation in the plane of the supplied figure
\end{definition}
\begin{proof}
  This is the declared model definition of Rotation Sense.
\end{proof}
\begin{definition}[Rotation Axis Location]
  \label{def:physics:phyx-mini-0306:phyxminiproblems-problemphyxmini0306-rotationaxislocation}
  \lean{PhyXMiniProblems.ProblemPhyXMini0306.RotationAxisLocation}
  Location of the rotation axis relative to the loop
\end{definition}
\begin{proof}
  This is the declared model definition of Rotation Axis Location.
\end{proof}
\begin{definition}[Gravity Regime]
  \label{def:physics:phyx-mini-0306:phyxminiproblems-problemphyxmini0306-gravityregime}
  \lean{PhyXMiniProblems.ProblemPhyXMini0306.GravityRegime}
  Whether gravity contributes appreciably to the string tension
\end{definition}
\begin{proof}
  This is the declared model definition of Gravity Regime.
\end{proof}
\begin{definition}[String Excitation]
  \label{def:physics:phyx-mini-0306:phyxminiproblems-problemphyxmini0306-stringexcitation}
  \lean{PhyXMiniProblems.ProblemPhyXMini0306.StringExcitation}
  How the string is excited in the scenario
\end{definition}
\begin{proof}
  This is the declared model definition of String Excitation.
\end{proof}
\begin{definition}[Wave Motion Kind]
  \label{def:physics:phyx-mini-0306:phyxminiproblems-problemphyxmini0306-wavemotionkind}
  \lean{PhyXMiniProblems.ProblemPhyXMini0306.WaveMotionKind}
  The kind of disturbance whose propagation speed is requested
\end{definition}
\begin{proof}
  This is the declared model definition of Wave Motion Kind.
\end{proof}
\begin{definition}[Circular Loop Figure]
  \label{def:physics:phyx-mini-0306:phyxminiproblems-problemphyxmini0306-circularloopfigure}
  \lean{PhyXMiniProblems.ProblemPhyXMini0306.CircularLoopFigure}
  \uses{def:physics:phyx-mini-0306:phyxminiproblems-problemphyxmini0306-rotationsense}
  The primary image contains one unlabelled circular outline and a curved clockwise arrow. It does not visibly mark the center point; rotation about the center comes from the prose rather than from an invented figure label
\end{definition}
\begin{proof}
  This is the declared model definition of Circular Loop Figure.
\end{proof}
\begin{definition}[Rotating Circular String Setup]
  \label{def:physics:phyx-mini-0306:phyxminiproblems-problemphyxmini0306-rotatingcircularstringsetup}
  \lean{PhyXMiniProblems.ProblemPhyXMini0306.RotatingCircularStringSetup}
  \uses{def:physics:phyx-mini-0306:phyxminiproblems-problemphyxmini0306-lengthquantity, def:physics:phyx-mini-0306:phyxminiproblems-problemphyxmini0306-massquantity, def:physics:phyx-mini-0306:phyxminiproblems-problemphyxmini0306-linearmassdensityquantity, def:physics:phyx-mini-0306:phyxminiproblems-problemphyxmini0306-accelerationquantity, def:physics:phyx-mini-0306:phyxminiproblems-problemphyxmini0306-tensionquantity, def:physics:phyx-mini-0306:phyxminiproblems-problemphyxmini0306-speedquantity, def:physics:phyx-mini-0306:phyxminiproblems-problemphyxmini0306-rotationsense, def:physics:phyx-mini-0306:phyxminiproblems-problemphyxmini0306-rotationaxislocation, def:physics:phyx-mini-0306:phyxminiproblems-problemphyxmini0306-gravityregime, def:physics:phyx-mini-0306:phyxminiproblems-problemphyxmini0306-stringexcitation, def:physics:phyx-mini-0306:phyxminiproblems-problemphyxmini0306-wavemotionkind, def:physics:phyx-mini-0306:phyxminiproblems-problemphyxmini0306-circularloopfigure}
  Physical quantities of the ideal rotating loop. arcMassPerRadian is the mass of an infinitesimal arc divided by its dimensionless central angle. inwardTensionResultantPerRadian is the corresponding limiting inward force from the endpoint tensions. Keeping these quantities explicit separates the local circular dynamics from the target wave speed
\end{definition}
\begin{proof}
  This is the declared model definition of Rotating Circular String Setup.
\end{proof}
\begin{definition}[Matches Problem Scenario]
  \label{def:physics:phyx-mini-0306:phyxminiproblems-problemphyxmini0306-matchesproblemscenario}
  \lean{PhyXMiniProblems.ProblemPhyXMini0306.MatchesProblemScenario}
  \uses{def:physics:phyx-mini-0306:phyxminiproblems-problemphyxmini0306-rotatingcircularstringsetup}
  Qualitative physical information stated in the prose
\end{definition}
\begin{proof}
  This is the declared model definition of Matches Problem Scenario.
\end{proof}
\begin{definition}[Matches Supplied Figure]
  \label{def:physics:phyx-mini-0306:phyxminiproblems-problemphyxmini0306-matchessuppliedfigure}
  \lean{PhyXMiniProblems.ProblemPhyXMini0306.MatchesSuppliedFigure}
  \uses{def:physics:phyx-mini-0306:phyxminiproblems-problemphyxmini0306-rotatingcircularstringsetup}
  Primary-image evidence: a single circular outline and a curved clockwise arrow are visible, with no marked or text-labelled center. This predicate contains no wave-speed value
\end{definition}
\begin{proof}
  This is the declared model definition of Matches Supplied Figure.
\end{proof}
\begin{definition}[Matches Problem Readouts]
  \label{def:physics:phyx-mini-0306:phyxminiproblems-problemphyxmini0306-matchesproblemreadouts}
  \lean{PhyXMiniProblems.ProblemPhyXMini0306.MatchesProblemReadouts}
  \uses{def:physics:phyx-mini-0306:phyxminiproblems-problemphyxmini0306-lengthincentimeters, def:physics:phyx-mini-0306:phyxminiproblems-problemphyxmini0306-speedincentimeterspersecond, def:physics:phyx-mini-0306:phyxminiproblems-problemphyxmini0306-rotatingcircularstringsetup}
  The two numerical readouts supplied by the problem statement
\end{definition}
\begin{proof}
  This is the declared model definition of Matches Problem Readouts.
\end{proof}
\begin{definition}[Has Physical Parameters]
  \label{def:physics:phyx-mini-0306:phyxminiproblems-problemphyxmini0306-hasphysicalparameters}
  \lean{PhyXMiniProblems.ProblemPhyXMini0306.HasPhysicalParameters}
  \uses{def:physics:phyx-mini-0306:phyxminiproblems-problemphyxmini0306-lengthreadout, def:physics:phyx-mini-0306:phyxminiproblems-problemphyxmini0306-linearmassdensityreadout, def:physics:phyx-mini-0306:phyxminiproblems-problemphyxmini0306-speedreadout, def:physics:phyx-mini-0306:phyxminiproblems-problemphyxmini0306-rotatingcircularstringsetup}
  Strict positivity conditions needed for the nondegenerate circular motion and for cancellation in the wave-speed comparison. The remaining quantities are already nonnegative because their underlying scalar type is NNReal
\end{definition}
\begin{proof}
  This is the declared model definition of Has Physical Parameters.
\end{proof}
\begin{definition}[Satisfies Rotating Loop Dynamics]
  \label{def:physics:phyx-mini-0306:phyxminiproblems-problemphyxmini0306-satisfiesrotatingloopdynamics}
  \lean{PhyXMiniProblems.ProblemPhyXMini0306.SatisfiesRotatingLoopDynamics}
  \uses{def:physics:phyx-mini-0306:phyxminiproblems-problemphyxmini0306-lengthreadout, def:physics:phyx-mini-0306:phyxminiproblems-problemphyxmini0306-massreadout, def:physics:phyx-mini-0306:phyxminiproblems-problemphyxmini0306-linearmassdensityreadout, def:physics:phyx-mini-0306:phyxminiproblems-problemphyxmini0306-accelerationreadout, def:physics:phyx-mini-0306:phyxminiproblems-problemphyxmini0306-tensionreadout, def:physics:phyx-mini-0306:phyxminiproblems-problemphyxmini0306-speedreadout, def:physics:phyx-mini-0306:phyxminiproblems-problemphyxmini0306-rotatingcircularstringsetup}
  Local dynamics of a uniform string element in steady circular motion: arc mass per radian is μ r; centripetal acceleration is v² / r; the limiting inward resultant of the two endpoint tensions, per radian, is the tension magnitude T; Newton's second law equates that resultant with mass times acceleration. These laws imply T = μ v²; they do not mention the requested wave speed
\end{definition}
\begin{proof}
  This is the declared model definition of Satisfies Rotating Loop Dynamics.
\end{proof}
\begin{definition}[Satisfies Transverse String Wave Law]
  \label{def:physics:phyx-mini-0306:phyxminiproblems-problemphyxmini0306-satisfiestransversestringwavelaw}
  \lean{PhyXMiniProblems.ProblemPhyXMini0306.SatisfiesTransverseStringWaveLaw}
  \uses{def:physics:phyx-mini-0306:phyxminiproblems-problemphyxmini0306-linearmassdensityreadout, def:physics:phyx-mini-0306:phyxminiproblems-problemphyxmini0306-tensionreadout, def:physics:phyx-mini-0306:phyxminiproblems-problemphyxmini0306-speedreadout, def:physics:phyx-mini-0306:phyxminiproblems-problemphyxmini0306-rotatingcircularstringsetup}
  The standard constitutive law for transverse waves on an ideal uniform string, T = μ c², expressed in every compatible system of base units. This is a governing law, not the conclusion c = v and not a numerical answer
\end{definition}
\begin{proof}
  This is the declared model definition of Satisfies Transverse String Wave Law.
\end{proof}
\begin{lemma}[tension eq linear Density mul tangential Speed sq]
  \label{lem:physics:phyx-mini-0306:phyxminiproblems-problemphyxmini0306-tension-eq-lineardensity-mul-tangentialspeed-sq}
  \lean{PhyXMiniProblems.ProblemPhyXMini0306.tension_eq_linearDensity_mul_tangentialSpeed_sq}
  \uses{def:physics:phyx-mini-0306:phyxminiproblems-problemphyxmini0306-linearmassdensityreadout, def:physics:phyx-mini-0306:phyxminiproblems-problemphyxmini0306-tensionreadout, def:physics:phyx-mini-0306:phyxminiproblems-problemphyxmini0306-speedreadout, def:physics:phyx-mini-0306:phyxminiproblems-problemphyxmini0306-rotatingcircularstringsetup, def:physics:phyx-mini-0306:phyxminiproblems-problemphyxmini0306-hasphysicalparameters, def:physics:phyx-mini-0306:phyxminiproblems-problemphyxmini0306-satisfiesrotatingloopdynamics}
  The rotating-loop laws imply the intermediate tension relation T = μ v². It is derived rather than included as a premise
\end{lemma}
\begin{proof}
  The conclusion follows from the governing relations and figure readouts named in the statement of tension eq linear Density mul tangential Speed sq.
\end{proof}
\begin{definition}[Answer Choice]
  \label{def:physics:phyx-mini-0306:phyxminiproblems-problemphyxmini0306-answerchoice}
  \lean{PhyXMiniProblems.ProblemPhyXMini0306.AnswerChoice}
  Labels of the four speed choices printed with the problem
\end{definition}
\begin{proof}
  This is the declared model definition of Answer Choice.
\end{proof}
\begin{definition}[speed Centimeters Per Second]
  \label{def:physics:phyx-mini-0306:phyxminiproblems-problemphyxmini0306-answerchoice-speedcentimeterspersecond}
  \lean{PhyXMiniProblems.ProblemPhyXMini0306.AnswerChoice.speedCentimetersPerSecond}
  \uses{def:physics:phyx-mini-0306:phyxminiproblems-problemphyxmini0306-answerchoice}
  Speed in centimeters per second printed beside each answer label
\end{definition}
\begin{proof}
  This is the declared model definition of speed Centimeters Per Second.
\end{proof}
\begin{definition}[recorded Dataset Answer]
  \label{def:physics:phyx-mini-0306:phyxminiproblems-problemphyxmini0306-recordeddatasetanswer}
  \lean{PhyXMiniProblems.ProblemPhyXMini0306.recordedDatasetAnswer}
  \uses{def:physics:phyx-mini-0306:phyxminiproblems-problemphyxmini0306-answerchoice}
  Answer label recorded by the source dataset
\end{definition}
\begin{proof}
  This is the declared model definition of recorded Dataset Answer.
\end{proof}
% --- Archon declaration topology end ---
% --- Archon physics formalization source end ---
